\documentclass{article}
\usepackage[UTF8]{ctex}
\usepackage{amsmath,amssymb,amsfonts}
\usepackage{enumitem}
\usepackage{geometry}
\geometry{a4paper, margin=1in}
% === 错题本排版优化（不参与题目标识） ===
\setlength{\parindent}{0pt}
\setlength{\parskip}{0.35em plus 0.10em minus 0.08em}
\setlength{\emergencystretch}{3em}
\setlist{itemsep=0.25em, topsep=0.40em, parsep=0pt}
\allowdisplaybreaks[2]
% === 错题本排版优化结束 ===
\begin{document}
\section*{上海大学2024年数学分析真题}
\section*{45．上海大学2024年数学分析试卷}
\begin{itemize}

\subsection{题目1}
  \item[1．] （10 分）叙述函数极限 \(\lim _{x \rightarrow x_{0}^{-}} f(x)=A \in \mathbb{R}\) 的＂\(\varepsilon-\delta\)＂定义．另外，利用函数极限的定义来证明
\[
\lim _{x \rightarrow 2^{-}} \frac{x^{2}-4}{x}=0 .
\]

\subsection{题目2}
  \item[2．] （10 分）假设数列 \(\left\{a_{n}\right\}\) 满足 \(\lim _{n \rightarrow \infty}\left|\frac{a_{n}}{a_{n+1}}\right|=\lambda>1\) ．证明： \(\lim _{n \rightarrow \infty} a_{n}=0\) ．
\end{itemize}

\subsection{题目3}
3．（10 分）计算极限 \(\lim _{n \rightarrow \infty}\left(\cos \frac{1}{\sqrt{n}}\right)^{n}\) 与 \(\lim _{x \rightarrow 0} \frac{\sqrt{1+\tan x}-\sqrt{1+\sin x}}{x \ln (1+x)-x^{2}}\) ．

\subsection{题目4}
4．（10 分）讨论函数 \(y=\sin (\sin x)\) 与 \(y=x \cos x\) 在 \([0,+\infty)\) 上的一致连续性并给出相应证明．

\subsection{题目5}
5．（10 分）设函数 \(f(x) \in C[0,1]\) ，且 \(4 \int_{0}^{1} f(u) \mathrm{d} u=1\) ，证明：存在 \(\theta \in[0,1]\) ，使得 \(f(\theta)=\theta^{3}\) ．

\subsection{题目6}
6．（10 分）求幂级数 \(\sum_{n=0}^{\infty}(-1)^{n} \frac{x^{n}}{n+1}\) 的收敛半径，收敛域与和函数，并给出 \(\sum_{n=0}^{\infty}(-1)^{n} \frac{1}{(n+1) 2^{n}}\) 的和．

\subsection{题目7}
7．（10 分）求空间曲线 \(\left\{\begin{array}{l}x^{2}+y^{2}=2 z ; \\ x^{2}+y^{2}+x y=1 .\end{array}\right.\) 上的点到 \(x O y\) 平面的最大与最小距离．

\subsection{题目8}
8．（10 分）设 \(A, B\) 为有界数集且 \(A \cap B \neq \emptyset\) 。证明： \(\sup (A \cap B) \leq \min \{\sup A, \sup B\}\) ，并给出等号不成立的例子。

\subsection{题目9}
9．（10 分）设 \(f(x) \geq 0\) 在 \([0,+\infty)\) 上可微，\(f(0)=0,\left|f^{\prime}(x)\right| \leq f(x)\) ．证明：在 \([0,+\infty)\) 上 \(f(x) \equiv 0\) ．

\subsection{题目10}
10．（10 分）讨论函数 \(z=f(x, y)=\left\{\begin{array}{ll}\frac{x y^{3}}{x^{2}+y^{4}}, & x^{2}+y^{2} \neq 0 ; \\ 0, & x^{2}+y^{2}=0 .\end{array}\right.\) 在原点 \((0,0)\) 的连续性，可偏导性及可微性，并给出证明．

\subsection{题目11}
11．（10 分）讨论反常积分 \(\int_{2}^{+\infty} \frac{1}{x \ln ^{2} x} \mathrm{~d} x\) 与 \(\int_{1}^{+\infty} \frac{1}{x^{2} \sqrt{\ln x}} \mathrm{~d} x\) 的敛散性并给出相应证明．

\subsection{题目12}
12．（10 分）将函数 \(f(x)=1-x^{2}(0 \leq x \leq \pi)\) 展开成余弦级数，并求级数 \(\sum_{n=1}^{\infty} \frac{(-1)^{n-1}}{n^{2}}\) 的和．

\subsection{题目13}
13．（10 分）求第二类曲面积分

\[
\iint_{S} x^{2} \mathrm{~d} y \mathrm{~d} z+y^{2} \mathrm{~d} z \mathrm{~d} x+z^{2} \mathrm{~d} x \mathrm{~d} y
\]

其中 \(S\) 是球面 \((x-a)^{2}+(y-b)^{2}+(z-c)^{2}=R^{2}\) ，并取外侧．

\subsection{题目14}
14．（10 分）计算曲线积分

\[
I=\int_{L}\left(e^{x} \sin y-2 x-2 y\right) \mathrm{d} x+\left(e^{x} \cos y-x\right) \mathrm{d} y
\]

其中 \(L\) 为从 \(O(0,0)\) 沿曲线 \(y=\sqrt{2 x-x^{2}}\) 到 \(A(2,0)\) 的一段有向弧．

\subsection{题目15}
15．（10 分）假设 \(f(x)\) 在 \((-\infty,-1]\) 上连续且 \(\int_{-\infty}^{-1} f(x) \mathrm{d} x\) 收敛，证明：存在数列 \(\left\{x_{n}\right\} \subset(-\infty,-1]\) ，使得流足 \(\lim _{n \rightarrow \infty} x_{n}=-\infty\) 且 \(\lim _{n \rightarrow \infty} f\left(x_{n}\right)=0\) 。另外，假设函数 \(g(x):(0,+\infty) \rightarrow(0,+\infty)\) 可导，证明：存在趋于正无穷大的正数列 \(\left\{x_{n}\right\}\) ，使得 \(g^{\prime}\left(x_{n}\right)<g\left(2 x_{n}\right), n=1,2, \cdots\) ．
\end{document}
